% Shared by the guide and tutorial; compile from their respective directories.
\section*{How to use Figaro today}
Figaro separates a probabilistic model from the algorithms that reason about it.
The modernized development line uses JDK 17, Scala 3.9.0 LTS and sbt 2.0.8.
It retains the \texttt{com.cra.figaro} packages, but its Scala 3 artifact is not
binary compatible with old Scala 2 libraries. This is a development baseline,
not an announcement of a published stable release.

\section*{Quick start}
Install JDK 17, Git and an sbt runner. Clone the repository and run its checked
example from the repository root:
\begingroup\scriptsize
\begin{verbatim}
git clone https://github.com/mattwilkinsphoto/figaro.git
cd figaro
sbt "examples / Compile / runMain com.cra.figaro.example.documentation.QuickStart"
\end{verbatim}
\endgroup
The checkout pins sbt and Scala; a separate Scala installation is unnecessary.
For dependency coordinates and compiled-library integration, use the current
\href{https://github.com/mattwilkinsphoto/figaro/blob/main/README.md}{README}
and \href{https://github.com/mattwilkinsphoto/figaro/blob/main/CONSUMER_BOUNDARY.md}{consumer contract}.
Do not reuse FigaroWork installers, old Scala 2 coordinates, or instructions to
put a library JAR on the system PATH.

\section*{Keep the foundations; follow the bridge for changes}
The original tutorial explains representation, evidence, relational models,
inference, filtering, decisions, learning and extensions in depth. Its chapters
are retained as historical teaching material, not asserted to be fully ported
Scala 3 listings. Original notices and authorship remain unchanged.
The \href{https://github.com/mattwilkinsphoto/figaro/blob/main/docs/DOCUMENTATION_MIGRATION.md}{chapter-by-chapter bridge}
connects that material to current guidance without reproducing the book.
The \href{https://github.com/mattwilkinsphoto/figaro/blob/main/docs/TUTORIAL.md}{classic walkthrough}
provides checked Scala 3 Constant, Select and Burglary examples.
Use the \href{https://github.com/mattwilkinsphoto/figaro/blob/main/docs/MIGRATION.md}{migration guide}
for changed syntax, dependencies and accepted workarounds, and the
\href{https://github.com/mattwilkinsphoto/figaro/blob/main/docs/API_GUIDE.md}{API guide}
for current contracts. Generate the searchable API using
\texttt{sbt "figaro / Compile / doc"}; the old Scala 2 API archive is obsolete.

\section*{Optional performance and stopping features}
Parallel importance sampling and multi-chain MCMC use isolated model instances
and bounded workers. They are explicit choices, not automatic acceleration of
existing programs or permission to share a mutable universe across threads.
Continuous-vector samplers operate on an explicit log density rather than an
arbitrary Figaro graph. Diagnostics, proposal calibration and stopping criteria
have distinct assumptions and budgets; passing a precision rule does not prove
that a chain has explored every important mode.
See the \href{https://github.com/mattwilkinsphoto/figaro/blob/main/docs/PARALLEL_PERFORMANCE.md}{parallel guide},
\href{https://github.com/mattwilkinsphoto/figaro/blob/main/docs/MULTI_CHAIN_MCMC.md}{multi-chain guide}
and \href{https://github.com/mattwilkinsphoto/figaro/blob/main/docs/STOPPING_CRITERIA.md}{stopping guide}.
Measured gains depend on workload and include overhead; there is no universal
speedup promise. Some historical tests remain flaky or failing, and OSGi and
unrestricted shared-graph concurrency remain outside the validated boundary.

\section*{Modernization credits}
Matthew Wilkins: modernization direction, requirements, review and validation.
Codex (OpenAI): AI-assisted implementation, testing, performance analysis and
documentation, working with Matthew Wilkins. These supplement the original
Figaro contributors; see \texttt{FigaroAttributions.txt}. Copyright and license
notices are not replaced by these credits.
